\documentclass[12pt]{amsproc}
\usepackage{mynotes}

\usepackage{enumitem}

% To list statements of a theorem
\newlist{statements}{enumerate}{1}
\setlist[statements]{
    label=\textbf{\arabic*)},    
    ref=\textbf{\arabic*)}               
}
% Macro to write implications
\newcommand{\imp}[2]{\ref{#1} $\implies$ \ref{#2}}

\newcommand{\cU}{\mathcal{U}}
\renewcommand{\P}{\mathbb{P}}

\newcommand{\upnil}{\operatorname{Nil}}

\renewcommand{\course}{Polynomials}

\title{\course}
\author{Leandro Vendramin}
\address{Department of Mathematics and Data
Science, Vrije Universiteit Brussel, Pleinlaan 2, 1050 Brussel}
\email{Leandro.Vendramin@vub.be}
\thanks{}
\date{}

\begin{document}

\begin{abstract}
\end{abstract}

\maketitle

\tableofcontents

\section{Introduction}

\section{}

Let $A$ be a commutative ring with unity $1\ne 0$. If $R$ is a subring 
of $A$ and $r\in R$, we denote by $R[r]$ the smallest subring 
of $A$ containing $R$ and $r$, that is, the intersection of all subrings of $A$ containing $R$ and $r$.
This sounds a bit abstract, but it is not. It is nothing but
the set of polynomial expressions 
in $r$ with coefficients in $R$, that is,
\[
R[r] = \{ a_0 + a_1 r + \cdots + a_n r^n : n\ge 0, a_i\in R \}.
\]

\begin{example}
Let $R$ be a subring of $A$ and let $r\in R$. Then $R[r]=R$. In particular, 
$R[1]=R[0]=R$.
\end{example}

\begin{example}\
    \begin{enumerate}
        \item $\Z[\sqrt{2}]=\{ a + b \sqrt{2} : a,b\in \Z \}$.
        \item $\Q[\sqrt{3}]=\{ a + b \sqrt{3} : a,b\in \Q \}$.
        \item $\Q[\sqrt[3]{5}]=\{ a + b \sqrt[3]{5} + c \sqrt[3]{25} : a,b,c\in \Q \}$.
    \end{enumerate}
\end{example}

\begin{theorem}
    Let $R$ be a subring of $A$ and let $a\in A$. Then
    $R[a]$ is a subring of $A$.
\end{theorem}

\begin{proof}

\end{proof}

\begin{exercise}
    There are no ring homomorphisms $\Z[\sqrt{2}] \to \Z[\sqrt{3}]$. In particular, 
    the rings $\Z[\sqrt{2}]$ and $\Z[\sqrt{3}]$ are not isomorphic. 
\end{exercise}


\begin{definition}
    Let $R$ be a subring of $A$ and let $a\in A$. 
    We say that $a$ is \emph{transcendental over $R$} 
    if 
    \[
    r_0+r_1a+\cdots+r_na^n=0\implies r_0=r_1=\cdots=r_n=0
    \]
    for all $r_0,\dots,r_n\in R$. If $a$ is not transcendental over $R$, we say that $a$ is \emph{algebraic over $R$}.
\end{definition}

\begin{exercise}
    Let $a\in A$ be transcendental over $R$... Prove that 
    $f(a)=g(a)$ in $R[a]$ if and only if $f=g$ in $R[x]$. 
\end{exercise}

Note that $\sqrt{2}$ is algebraic over $\Z$ because
there is no unique algebraic expression for it with integer coefficients. For example,
\[
\sqrt{2}=-2+\sqrt{2}+(\sqrt{2})^2.
\]

\begin{theorem}
Let $x,y\in R$ be transcendental over $R$. Then $R[x]\simeq R[y]$ as rings.
\end{theorem}

\begin{proof}

\end{proof}

In the same way, we can prove that if $R$ is a subring of both $A$ and $B$, 
and $a\in A$ and $b\in B$ are transcendental over $R$, then $R[a]\simeq R[b]$ as rings.

\begin{definition}
    Let $A$ be a commutative ring with unity $1\ne 0$, and $R$ be a subring of $A$. Let 
    $x\in A$ be transcendental over $R$. The \emph{polynomial ring} $R[X]$ 
    in the indeterminate $X$ 
    with coefficients in $R$ is the ring of polynomial expressions 
    \[
    r_0+r_1X+\cdots+r_nX^n
    \]
    together with a ring isomorphism $f\colon R[X]\to R[x]$ such that $f(X)=x$. 
\end{definition}

Note that the polynomial ring $R[X]$ is unique up to a unique isomorphism. If $y$ is another trascendental element 
over $R$, then $R[X]\simeq R[x]\simeq R[y]$ as rings.

\framebox{existence} Very technical, see the appendix. 

We often write 
\[
f(X)=\sum_{i=0}^n r_iX^i=r_0+r_1X+r_2X^2+\cdots+r_nX^n,
\]
where, by convention, $X^0=1$. 

\begin{definition}
    The \emph{degree} of a non-zero polynomial...

\end{definition}

The degree is well-defined because...

The zero polynomial $0$ is the only polynomial that does not have a degree.

\begin{example}
    The polynomial $f(X)=\sum_{i=0}^n r_iX^i\in R[X]$ has degree $n$ if and only if $r_n\neq 0$. For example, 
    $3+X+0X^6$ has degree one. 
\end{example}

\begin{definition}
    Let $f(X)=\sum_{i=0}^n r_iX^i\in R[X]$ be a polynomial of degree $n$. 
    We say that $f(X)$ is \emph{monic} if $r_n=1$. The \emph{leading coefficient} of $f(X)$ is $r_n$, 
    and the \emph{constant term} of $f(X)$ is $r_0$.
\end{definition}

\begin{exercise}
    Let $f(X)=...$ and $g(X)=...$. Compute...
\end{exercise}

\begin{example}
    Find polynomials $q(X),r(X)\in \Z[X]$ such that...
\end{example}

\begin{exercise}
    Find all polynomials $f(X)\in \Z_2[X]$ of degree at most $3$. How many are there? How many are monic?
\end{exercise}

\begin{theorem}
Let $f(X),g(X)\in R[X]$ be non-zero polynomials. Then 
\[ 
\deg(f(X)g(X))\leq \deg(f(X))+\deg(g(X)).
\] 
If $R$ is an integral domain, 
then $\deg(f(X)g(X))=\deg(f(X))+\deg(g(X))$.
\end{theorem}

\begin{proof}

\end{proof}

\begin{corollary}
    $R[X]$ is an integral domain if and only if $R$ is an integral domain.
\end{corollary}

\begin{proof}

\end{proof}

\begin{example}
If $R$ is an integral domain and $f(X)\in R[X]$ 
has degree $n$, then $f(X)^m$ has degree $nm$ for all $m\ge 1$. In particular, $f(X)^m\neq 0$ for all $m\ge 1$.
\end{example} 

\begin{example}
    Let us prove that there are no polynomials $f(X),g(X)\in \Z_2[X]$ such that 
    $f(X)^2+g(X)^2=X$. Indeed, if $f(X)=\sum_{i=0}^n a_iX^i$ and $g(X)=\sum_{i=0}^m b_iX^i$, then
    ...
\end{example}

\begin{exercise}
    Let $K$ be a field. Prove that there no polynomials $f(X),g(X)\in K[X]$ such that 
    $f(X)^2+Xg(X)^2=0$.
\end{exercise}

\section{Units} 

\begin{proposition}
    Let $R$ be an integral domain and $f(X)\in R[X]$. Then $f(X)$ is a unit in $R[X]$ if and only if $f(X)$ is a unit in $R$.
\end{proposition}

\begin{proof}

\end{proof}

\begin{example}
    $(1+2X)(1+2X)=1$ in $\Z_4[X]$, but $1+2X$ has degree one.
\end{example}

\begin{exercise}
    Compute the units of $(\Z/4\Z)[X]$ and $(\Z/5\Z)[X]$.
\end{exercise}

\section{Application: binomial coefficients}

This section is optional, but shows a nice application of polynomials.

\section{Divisibility}

\section{Appendix: Existence of polynomial rings}

This section is rather technical and can be skipped. It is included for completeness.

\bibliographystyle{abbrv}
\bibliography{refs}

\end{document}
